\documentclass[11pt,a4paper]{amsart}
\usepackage[margin=30mm]{geometry} 
\usepackage{amssymb}
\usepackage{amsthm}
\usepackage{url}
\usepackage{bbm} 
\usepackage{listings} 
\newcommand{\Z}{{\mathbb{Z}}}
\newcommand{\R}{{\mathbb{R}}}
\newcommand{\Q}{{\mathbb{Q}}}
\newcommand{\N}{{\mathbb{N}}}
\newcommand{\F}{\mathbb{F}}
\newcommand{\C}{\mathbb{C}}
\newcommand{\pow}{\operatorname{pow}}

\newtheorem{thm}{Theorem} 
\newtheorem{rem}{Remark} 
\newtheorem{lem}{Lemma} 


\title[Residue class rings and the CRT]{Residue class rings, Chinese Remainder Theorem, invertibility, fields from the quotient construction\\ (Lecture notes)}
\author{G. Averkov} 
\begin{document} 
	\maketitle
	\vspace{-2em}
	\begin{center} 
		\today 
	\end{center} 
	
\section{Residue class rings} 

The rings in this section are all unitary commutative rings. It will be convenient to give up the condition $0_R \ne 1_R$ in the definition of the rings, that is, the rings that we consider here maybe of the form $R = \{0\}$. This is called a zero or a trivial ring. Normally, the trivial ring is just a trivial case in the discussion that follow. 

\subsection{Ring homomorphism, ring isomorphisms and isomorphic rings} For rings $R,S$ a map $\phi : R \to S$ is called a \emph{ring homomorphism} if respects the structure of the ring, meaning that the equalities $\phi(u+v) = \phi(u)+\phi(v)$ and $\phi(u \cdot v) = \phi(u) \cdot \phi(v)$ hold for all $u,v \in R$ and, apart from that, one has the equalities $\phi(0) = 0$, $\phi(1) = 1$. 

We remark that the condition $\phi(0)=0$ required here is redundant as $\phi(0) = \phi(0+0)=\phi(0) + \phi(0)$ gives $\phi(0) = 0$. The reason we want to keep in the definition is the fact that all the special constants of the structure and the special operations of the structure get preserved. This is directly seen from the definition we give.

A ring homomorphism is called a ring isomorphism if it is bijective. Two rings $S$ and $T$ for which the exists a \emph{ring isomorphism} $S \to T$ are called \emph{isomorphic} and one uses the notation $R \cong S$ in this case. Informally, isomorphic rings are equal up to a `translation' of notation. 

The \emph{kernel} of  a ring homomorphism $\phi : R \to S$ is defined by 
\[
	\ker(\phi) := \{ r \in R \colon \phi(r) = 0 \}. 
\]

\subsection{Monomorphisms and epimorphisms} There are the following names that are also used: injective homomorphisms are called \emph{monomorphisms} and surjective homomorphisms are called \emph{epimorphisms}. However, the usage of these terms would expand the vocabulary even further, so that we prefer just to say injective  homomorphism and surjective homomorphism. 

\subsection{Examples of ring homomorphisms} Here are some examples related for the rings $k[x_1,\ldots,x_n]$ with $n \in \{1,2\}$, where we choose the coefficient field to be $k = \R$ to make the presentation geometric in the truly visual sense. 
\begin{enumerate} 
	\item $\R[x,y] \to \R[x]$, $f(x,y) \mapsto f(x,x^2)$. The kernel of this homomorphism is the set of all polynomials that vanish on the parabola $y = x^2$. 
	\item $\R[x] \to \R[x]$, $f(x) \mapsto f(x-1)$ is a ring isomorphism.  
	\item $\R[x] \to \R[x]$, $f(x) \mapsto f(x^2)$ injective but  not surjective.  
\end{enumerate} 

\subsection{Kernel is an ideal} The kernel of $\phi : R \to S$ is an ideal of the ring $R$. In fact $0 \in \ker (\phi)$ as $\phi(0) = 0$. If $u,v \in \ker (\phi)$, then $u+v \in \ker (\phi)$ and if $u \in R$ and $\phi(v) = 0$, then $\phi(u v) = \phi(u) \phi(v) = \phi(u) 0 = 0$. 

\subsection{On injectivity} A ring homomorphism $\phi : R \to S$ is injective if and only if $\ker (\phi) = \{0\}$. The necessity is clear. For the sufficiency, take $u,v \in R$ with $\phi(u) = \phi(v)$. This implies $u-v \in \ker (\phi ) = \{0\}$ giving $u=v$. 

Thus, showing that a ring homomorphism  $\phi : R \to S$ is a ring isomorphism boils down to checking that $\phi$ is surjective and that $\ker (\phi) = \{0\}$. 

\subsection{Congruence modulo an ideal and modulo an element of the ring}  Elements $u,v \in R$ are said to be \emph{congruent modulo} an ideal $I \subseteq R$, written as $u \equiv v \mod I$, if $u - v \in I$. Congruence is an equivalence relation (it is reflexive, symmetric and transitive). The equivalence classes with respect to the congruence relation are called \emph{residue classes} or \emph{cosets} modulo the ideal $I$. 

Given an element  $f$ of a ring  $R$, we can consider the congruence relation and the cosets modulo the principal ideal $f R$ generated by $f$. In this case, we speak about the congruence and the cosets modulo $f$. That is, the relation $u \equiv v \mod f$ means $u - v \in f R$ and cosets modulo $f$ are cosets modulo the ideal $fR$.

\subsection{The unit ideal} An ideal $I$ coincides with the whole ring $R$ if and only if $1 \in I$. The ideal $I=R$ is called the unit ideal. 

\subsection{Residue class ring} Let $R$ be a ring and let $I$ be an ideal in $R$. The set $r + I := \{r + u \colon u \in I\}$ is a \emph{residue class} (or a \emph{coset}) of $r$ modulo $I$. Another notation for it is $[r]_I$. Let $R / I$ be the set of all residue classes. On that set we introduce the addition and the multiplication by 
\begin{align*} 
[u]_I + [v]_I & := [u+v]_I
\\
	[u]_I \cdot [v]_I & := [u \cdot v]_I, 
\end{align*} 
For this relation to be well defined, one needs to verify that the result on the right-hand side in the definition of the two operations does not depend on a particular representative of the cosets of the two operands. That is, one needs to check the implications
\begin{align*} 
		& u - r \in I, v - s \in I  & & \Longrightarrow &  (u + v ) - (r+s) \in I &, \\
		& u - r \in I, v- s \in I & & \Longrightarrow & u v - r s \in I. 
\end{align*} 
Both implications are not hard to check. 

 The \emph{residue class ring} is the structure defined as  the set $R / I$ with the above two operations $+$ and $\cdot$ attached to it. This is in fact a ring in our sense and it is a non-zero ring if $I \ne R$. If $I = R$, then $ R / I$ consists of just one element. Other common names for the residue class ring are the \emph{factor ring} of $R$ modulo $I$ or the \emph{quotient ring} of $R$ modulo $I$. 

Given an element $f$ of the ring $R$, we consider the above terminology for the principal ideal $ I = fR$ generated by $f$. In that case one speaks about the residue class ring modulo $f$ and uses the simplified notation $R / f$ for the ring $R / f R$ as well as the notation $[r]_f$ for the coset $[r]_{fR}$. 

\subsection{Examples of residue class rings} 
There are two basic we keep in mind and give illustration on, which are $R = \Z$ and $R= k [x_1,\ldots,x_n]$. Remember also that $\Z$ and $k[x]$ (polynomial ring with one variable) are analogous in many respects. Here are some example of residue class rings. 
\begin{enumerate} 
	\item $\Z / 2$ is not only a ring but actually a field. This is the field $\F_2$ with $2$ elements. 
	\item $ \Z / 3$, too, is not just a ring but actually a field. It is the field $\F_3$ with $3$ elements. 
	\item $\Z / 6$ is a ring with six elements (so, a finite ring), but not a field. 
	\item $\R[x] / (x^2+1)$ is isomorphic to the field $\C$ of complex numbers. Because the coset $i=[x]_{x^2+1}$ satisfies $i^2 + 1 = 0$. Hence, $\R[x] = \R[i]$ and this is essentially $\C$. 
	\item $\R[x] / (x^2 - 1) $ is a ring, but not a field. 
	\item  $\R[x,y] / (y-2  x)$ is isomorphic to $\R[x].$ 
	\item $\R[x,y,z] / \langle x, y \rangle $ is isomorphic to $\R[z]$. 
	\item Coordinate ring $R[X]  = \R[x_1,\ldots,x_n]/ I(X)$ coordinate ring of the set $X$. That is polynomials that are equal on $X$ get identified.
	\item $\R[x,y] / \langle x^2 + y^2 - 1 \rangle \cong \R[\cos t, \sin t]$
	\item $\Q[x] / \langle x^2 - 2\rangle \cong \Q[\sqrt{2}]$. 
\end{enumerate} 

\subsection{Canonical homomorphism of residue class rings} The operation of taking the coset of a given element $r \mapsto [r]_I$ from $R \to R / I$ is a ring homomorphism. This called the canonical ring homomorphism of $R / I$. 

Informally, coming from $r$ to $[r]_I$, we forget the whole information about $r$ and only keep the information modulo $I$. The statement above draws attention to the fact that we can do the arithmetic with $+$ and $\cdot$ consistently when only the partial information, modulo $I$, is available. 


In the context of $R = \Z$ and $R = k[x]$, the information modulo $I$ is a remainder information, because for these two rings all ideals have the form $I = f R$ for some $f \in R$. 

\subsection{Canonical homomorphism of two quotient rings} More generally, if we have two ideals $I$ and $J$ in $R$ and the map $[r]_I \mapsto [r]_J$ from $R/ I$ to $R / J$ is well-defined and is a ring homomorphism from $R/ I \to R / J$, which we also call canonical. 

For the rings $R = \Z$ and $R = k[x]$, this setting is the setting where $I = f R$ and $J = g R$ with $g$ dividing $f$. So, the assertion here is that we can pass to the remainders modulo $f$ to remainders modulo $g$ with all of the $+$ and $\cdot$ arithmetic operations being preserved. 

\section{Chinese Remainder Theorem} 

Chinese Remainder theorem is a convenient tool that allows one to break down calculations in a complicated ring into independent calculations in several less complicated rings. In the context of finite rings, this means that we break down a ring of a large size into several rings of a smaller size. Of course, the informal description is only useful if we make it rigorous and provide necessary details. 

\subsection{Operations on ideals} Let $I, J$ be ideals in a ring $R$. Then we introduce the sum the intersection of the ideals $I$ and $J$ as 
\begin{align*} 
	I + J & := \{ f+ g \colon f \in I, g \in J\}
\end{align*} 
This is in fact an ideal. We introduce the product of the two ideals $I$ and $J$ as the ideal $I \cdot J$ generated by all products of the form $f \cdot g$ with $f \in I$ and $g \in J$. This is also an ideal. And we consider the intersection $I \cap J$ of the ideals $I$ and $J$, which is again an ideal. 

We also observe that the three ideals we have introduced satisfy the inclusions 
\[
	I \cdot J \subseteq I \cap J \subseteq I + J
\] 
With respect to the Chinese Remainder Theorem, our focus will be mainly on $I \cap J$ and $I + J$, but we can also create a connection ti $I \cdot J$. Note that the ideals $I \cdot J$ and $I \cap J$ are not always the same. For example, $2 \Z \cdot 4 \Z = 8 \Z$, while $2 \Z \cap 4 \Z = 4 \Z$. But the equality of $I \cdot J$ and $I \cap J$ is possible. For example, we have $2 \Z \cap 3 \Z = 2 \Z \cdot 3 \Z = 6 \Z$. 

\subsection{Sum and the product of ideals in terms of generators} 
If $I = \langle f_1,\ldots, f_s\}$ and $J = \langle g_1,\ldots, g_t\rangle$, 
then 
\begin{align*} 
	I +J & =  \langle f_1 ,\ldots, f_s, g_1,\ldots, g_s \rangle, 
	\\ I \cdot J & = \langle f_i g_j \colon i=1,\ldots,s, \ j =1,\ldots,t\rangle 
\end{align*} 
This also demonstrates that $I+J$ corresponds to the intersection and $I \cdot J$ corresponds to the union of varieties in the setting $R = k[x_1,\ldots,x_n]$. 

\subsection{Coprime ideals} We call ideals $I$ and $J$ of a ring $R$ coprime if $I +  J  = R$. 

\subsection{Coprime ideals in $k[x]$} In the ring $k[x]$ each ideal is principal, so $I = \langle f \rangle$ and $J = \langle g \rangle$ holds and $I + J = \langle f, g \rangle = \langle \gcd(f,g) \rangle$. It follows that $I$ and $J$ are coprime if and only if $1= \gcd(f,g)$, that is, if the generators of $I$ and $J$ are coprime. The same observation is true for the ring $R= \Z$ of integers. 

\subsection{Product ring} The product ring $R \times S$ for two rings $R$ and $S$ is defined as the Cartesian product of $R$ and $S$, i.e., the set of all pairs $(r,s)$ with $r \in R$ and $s \in S$, with the addition and multiplication defined naturally in the component-wise fashion as 
\begin{align*} 
	(r,s) + (u,v) & := (r+u,s+v)
	\\ 
	(r,s) \cdot (u,v) & := (r \cdot u, s \cdot v)
\end{align*} 
This means in particular, that evaluation of an expression in a ring $R \times S$ boils down to evaluation the same kind of expression independently in the ring $R$ and in the ring $S$ and then putting the results together. Think for example, about calculating $x^2 - x + 1$ for $x= (r,s) \in R \times S$. The result is $(r^2 - r + 1, s^2 - s+1)$, where the two expressions $r^2 - r+1$ and $s^2 - s+1$ are the same but have been evaluated in their own rings ($R$ and $S$, respectively). 

Clearly, we can and will also consider the product of $m$ rings $R_1, \ldots, R_m$, which we denote by $R_1 \times \cdots \times R_m$.

The product of rings is also sometimes called a \emph{direct product}. 

\subsection{Chinese Remainder Theorem -- the version $R \to R / I \times R / J$.} Let $I,J$ be co-prime ideals in $R$. Then the homomorphism $\phi : R \to R / I \times R / J$, given by 
\[
	\phi(x) = ( [x]_I, [x]_J)
\] is surjective. Furthermore, $\ker(\phi) = I \cap J$. 
\begin{proof} 
	Effectively, we need to show that system of congruences
	\begin{align*} 
			\begin{cases} 
					x \equiv y \mod I,
					\\ 
					x \equiv z \mod J
			\end{cases} 
	\end{align*} 
	has a solution $x \in I$ for any given $y, z \in R$. Since $R = I+J$, we can write $1$ as $1 = u + v$ for some $u \in I$ and $v \in J$. Clearly, $u \equiv 0 \mod I$ and $v \equiv 0 \mod J$. But apart from this, $u = 1 - v \equiv 1 \mod J$ and $v = 1 -u \equiv 1 \mod I$. So, let us summarize this as 
	\begin{align*}
			v& \equiv 1 \mod I  & v & \equiv 0 \mod J
			\\ 
			u & \equiv 0 \mod I  & u & \equiv 1 \mod J. 
	\end{align*} 
That means, so to say, that $v$ is turned on on $I$ and turned of on $J$, while for $u$ it is the other way around, it is turned off on $I$ and turned on on $J$. Thus, in the spirit of Lagrange interpolation, we can build up a solution to our system of congruences. Indeed, $x = y v + z u$ solves our system of congruences giving an $x$ with $\phi(x) = ( [y]_I, [z]_J)$, 
because $x \equiv y \cdot 1 + z \cdot 0 \mod I$ and $x \equiv y \cdot 0 + z \cdot 1 \mod J$. 
	
	The kernel is $\{ x \in R \colon \phi(x) = 0\} = \{ x \in \R \colon [x]_I = 0, [x]_J = 0\} = \{x \in R \colon x \in I , x \in J\} = I \cap J.$
	
	Let us remark that the analogy with the Lagrange interpolation is not far fetched here, as we can in fact interpret the Lagrange interpolation in terms of the Chinese Remainder Theorem. 
\end{proof} 

\subsection{Extended Euclidean Algorithm} 
As we see from the proof of the CRT, we should be interested in an algorithm that provides a description of $1$ as $1 = u + v$ with $u \in I$ and $v \in J$. 
For both rings $R = \Z$ and $R = k[x]$, which are principal ideal domains the a long division algorithm, we can extended the Euclidean Algorithm to the so-called extended Euclidean algorithm (EEA) to find such a description. So, assume we are given the two ideals $I$ and $J$. As our ring $R$ is a principal ideal domain, they are of the form $I = a R$ and $I = b R$. So, assume that the ideals are given in this form via $a,b \in R$. What we need to find is a representation of $\gcd(a,b)$ as $\gcd(a,b) = a c + b d$ with $c,d \in R$. The basic idea of the Euclidean algorithm is to iteratively deploy the reduction $(a,b) \mapsto (b, a - q b)$, where $q$ is the quotient of the long division of $a$ by $b$ (assuming, of course, that $b \ne 0$). Since $\gcd(a,b) = \gcd(b,a-qb)$, we don't loss the value of the $\gcd$ but reduce the problem to a problem for smaller elements of $R$. What we can do now is assigning the variable names to $a$ and $b$. Let us call these names $A$ and $B$. respectively. That is we start with the expressions $A = a$ and $B = b$, where $a,b$ are concrete values in $R$, while $A$ and $B$ are the variable names for these values. When we pass from $(a,b)$ to $(b,a-qb)$ we can now attach symbolic expressions to both $b$ and $a-qb$. The symbolic expression for $b$ is $B$, while the symbolic expression for $a-q b= A - q B$. Say, if we have $A = 23$ and $B=5$, then we have reduced the calculation of $\gcd(A=23, B = 5)$ the calculation of $\gcd(B=5, A - 4 B = 3)$. So, as we iterate on and on, we can keep generating new linear expressions in $A$ and $B$ out of the existing ones. Say, in the next step we would come from $\gcd(B=5,A - 4 B =3)$ to $\gcd(A - 4B =3, B - (A-4B) =2) = \gcd(A - 4 B =3, 5 B - A = 2)$. Eventually, when we terminate, we have an expression for the $\gcd(a,b)$ written as a linear combination of $A=a$ and $B=b$. This is exactly what the Euclidean algorithm does. On the implementation level, to keep an arbitrary expression of the form $r A + s B =  t$ one can just use the vector $(r,s,t)$ and do transformations of such vectors. This is exactly what happens in Gaussian elimination method. So, the EEA is a sort of Gaussian elimination procedure type of algorithm, in which one starts with a solution and searches of an equation for this solution and in which the underlying structure is not a field, but a ring $R$, in which a long division is available. 



\subsection{Taking a homomorphism modulo the kernel} Every ring homomorphism $\phi : R \to S$ gives rise to an injective ring homomorphism by taking the quotient ring modulo $I = \ker \phi$. The injective ring homomorphism is $\psi : R / I \to S$, which is given by $\psi( [x]_I ) = \phi(x)$. It is well defined, because if $[u]_I = [v]_I$, then $u - v \in \ker(\phi)$, meaning that $\phi(u-v) = 0$, which means that $\phi(u) = \phi(v)$. Consequently, it does not matter what representative of the coset one takes. 

It is clear that the maps $\phi$ and $\psi$ have the same image, that is, $\phi(R) = \psi(R / I)$. 

Using this kind of transition, we can reformulate the Chinese Remainder Theorem as follows. 



\subsection{Chinese Remainder Theorem in the form $R / (I \cap J) \to R / I \times R / J.$} Let $I$ and $J$ be coprime ideals in a ring $R$. Then the map $\psi : R / (I \cap J) \to \R / I \times R / J$, 
$\psi([x]_{I \cap J}) = ([x]_I, [x]_J)$ is a ring isomorphism. 

\begin{proof} 
We have shown that the map $\phi : R \to R / I \times R / J$ with $\phi(x) = ([x]_I, [x]_J)$ is surjective with the kernel $I \cap J$. 

 Hence, $\psi$, too, is surjective. But, $\psi$ is injective, because its kernel is $0$: $\psi([x]_{I\cap J} ) = 0$  means that $[x]_I = 0$ and $[x]_J =0$, which means that $x \in I \cap J$, which means that $[x]_{I \cap J} = 0$. 
\end{proof} 

\subsection{A lemma on coprimality} If in a ring $R$, an ideal $I_1$ is coprime with each of the ideals $I_2,\ldots,I_m$. Then it is coprime with the ideal $I_2 \cap \cdots \cap I_m$. 
The ideal $I_1$ is coprime with each of the ideals $I_2,\ldots,I_m$. Hence, we can 
write$1 = u_i + v_i$ for some $u_i \in I_1$ and some $u_j \in I_j$ when $j=2,\ldots,m$. Taking the product of all these $m-1$ identities we arrive at 
\[
1 = \prod_{i=2}^m (u_i + v_i) 
\]
When we expand brackets we generate $2^{m-1}$ terms. As $u_i$ belongs to $I_1$ and $I_1$ is an ideal, whenever we pick $u_i$ from some of the brackets, we will end up with a term that belongs to $I_1$. The only other term that we will have is the one where we pick $v_i$ from each of the brackets. This gives the term $v_2 \cdots v_m$ belonging to all of the ideals $I_2, \ldots, I_m$ as $v_j \in I_j$ for $j = 2,\ldots,m$. So, the expansion of the brackets brings us to the representation 
\[
1 = a_1 + b_1,
\]
$a_1 \in I_1$ and $b_ 1= v_2 \cdots v_m \in I_2 \cap \cdots \cap I_m$. By the above argument we show that the ideals $I_1$ and $I_2 \cap \cdots \cap I_m$ are coprime. Actually, we even show something stronger: the ideals $I_1$ and $I_2 \cdot \ldots \cdot I_m$ are coprime. 


\subsection{Chinese Remainder theorem in the form $R \to R / I_1 \times \cdots \times R / I_m$} Let $I_1,\ldots,I_m$ be ideals in a ring $R$ that are pairwise coprime, meaning that $R = I_i + I_j$ for $i \ne j$. Then the map $\phi : R \to R / I_1 \times \cdots \times R / I_m $ given by $\phi(x) = ([x]_{I_1} , \ldots, [x]_{I_m})$ is a surjective ring homomorphism. 

\begin{proof} View of the lemma find representations $1 = a_i + b_i$ with $a_i \in I_i$ and $b_i \in I_j$ whenever $i \ne j$. Hence, $b_i =  1 - a_i \equiv 1 \mod I_i$ and $b_i \equiv 0 \mod I_j$ when $j \ne 0$. That is, on the ideal $I_i$ the element $b_i$ is turned on and on all the other ideals it is turned off. Thus, in the spirit of Lagrange interpolation, when we 
solve a system of the congruences $r \equiv y_i \mod I_i$ for some given $y_i \in I_i$ and the unknown $r \in R$, we can provide the solution 
\[
	r = b_1 y_1 + \cdots + b_m y_m
\]
to this system. 

\end{proof} 

\subsection{Chinese Remainder Theorem, the version $R / (I_1 \cap \cdots \cap I_m) \to R / I_1 \times \cdots R / I_m$}  Let $I_1,\ldots,I_m$ be ideals in $R$, then the map $\psi : R / (I_1 \cap \cdots \cap I_m) \to R / I_1 \times \cdots \times R / I_m$ given by $\psi( [x]_{I_1 \cap \cdots \cap I_m} = ([x]_{I_1},\ldots, [x]_{I_m})$ is a ring isomorphism. 

\begin{proof} Just use the previous version and pass to $ R / \ker \phi$. 
\end{proof} 

\subsection{Application to the integer ring} For the ring $\Z$, let us consider the smallest possible example where $m > 2$. So let us take $ \Z / 30 \cong\Z /2 \times \Z / 3 \times \Z / 5 \Z$. The arguments above show us that we can find an explicit way to invert the isomorphism 
\[
	[r]_{30} \mapsto ([r]_2,[r]_3,[r]_5). 
\]
To this end we need to run the EEA on the inputs $(2, 3 \cdot 5), (3, 2 \cdot 5)$ and $(5, \cdot 2 \cdot 3)$. We don't really have to do that as the outcome of the EEA is easy to guess on such small instances. The outputs are the representations 
\begin{align*} 
	1 & = -14 + 15 & & \text{with} & -14 & \in 2 \Z & & \text{and} & 15 & \in 3 \Z \cap 5 \Z,
\\
	1 & = -9 + 10 & & \text{with} & -9 & \in 3 \Z & & \text{and} & 10 & \in 2 \Z \cap 5\Z,
\\
	1 & = -5 + 6 & & \text{with} & -5 & \in 5 \Z & & \text{and} & 6 & \in 2 \Z \cap 3 \Z. 
\end{align*} 
So, $15$, $10$ and $6$ are turned on, respectively, on the ideals $2 \Z, 3 \Z$ and $5 \Z$ and are turned off on the other ideals. Hence, the system of congruence
\[
	\begin{cases}
		r  \equiv a \mod 2 \Z, 
\\		r  \equiv b \mod 3 \Z,
\\		r  \equiv c \mod 5 \Z,
	\end{cases} 
\]
with given right hand sides $a,b,c \in \Z$
is solved when
\[
	r \equiv 15 a + 10 b + 6 c \mod 30 \Z
\]
The inverse isomorphism is thus
\[
	([a]_2, [b]_3, [c]_5) \mapsto [15 a + 10 b + 6 c]_{30}. 
\]

\subsection{Application of the CRT to $k[x]$} The examples we considered for the CRT can also be given for the ring $R = k[x]$ with the calculations of isomorphisms carried out using  the EEA. As a side remark, we also would like to mention that there are rings other than $\Z$ and $k[x]$, which long division can be defined. Such rings are called Euclidean domains. 

\subsection{The CRT and polynomial interpolation} We have seen that when we prove the CRT we use some kind of interpolation idea with taking the polynomial modulo an ideal being an analogue of evaluation. There is indeed a very concrete connection between interpolation an the CRT when the underlying ring $R$ is the polynomial ring $k[x]$. We restrict our attention to the case of the coefficient field $k = \R$ for the sake of concreteness and present an example that highlights the connection. 

We start with the discussion of the simplest case, when build quotient rings modulo a degree one polynomial. As constants do not really matter, we can bring this polynomial into the form $x-a$ with $a \in \R$. For $a \in \R$, the map $\R[x] / (x-a) \R[x] \to \R$ given by $[f (x)]_{x-a} \mapsto f(a)$ is a ring isomorphism.  Indeed, the remainder after division of $f$ by $x-a$ is exactly $f(a)$, if the quotient-remainder representation of $f(x)$ is $f(x) = q(x) (x-a)+ r$, then plugging in $x=a$, we see that $f(a) = r$, that is, the remainder, is exactly $f(a)$. Hence, all polynomials in the same coset modulo $x-a$ have the same value at $a$. This also means that the canonical homomorphism $\R[x] \to \R[x] / (x-a) \R[x]$ that sends $f$ to $[f]_{x-a}$ can be interpreted as the evaluation at $a$, which is a map $\R[x] \to \R$ that sends $f$ to the value $f(a)$. 

Now, assume we are dealing with a more complicated quotient ring modulo a polynomial $g(x)$ that gets factorized into degree-one polynomials, say $g(x) = x (x-1)(x-2)$. This polynomial has roots $0, 1, 2$ each of multiplicity one. It is straightforward to see that the polynomials in the same coset modulo $g(x)$ all have the same values on the root set $\{0,1,2\}$. The converse is also true, if two polynomials have the same values on $\{0,1,2\}$, then their difference has roots at $0, 1$ and $2$, which means that they are divisible by $g(x)$. Thus, we see that the map $[f]_g \to (f(0),f(1),f(2))$ from the ring $\R[x] / g(x) \R[x]$ to the ring $\R \times \R \times \R$ is a ring isomorphism. The problem of inverting this isomorphism is essentially the problem of polynomial interpolation. 

The isomorphy $\R[x] / g(x) \R[x]$ and $\R \times \R \times \R$ can also be seen via the CRT theorem. Clearly 
\[
	g(x) \R[x] = x \R[x] \cap (x-1) \R[x] \cap (x-2) \R[x],
\]
where any two ideals in the intersection on the right-hand side are coprime. So, by the CRT, 
\[
	\R[x] / g(x) \R[x] \cong  \R[x] / x \R[x] \times  \R[x] / (x-1) \R[x] \times \R[x] / (x-2) \R[x] 
\]
but we have seen  that $\R[x] / (x-a) \R[x] \cong \R$ via the evaluation $[f]_{x-a} \mapsto f(a)$. Hence 
\[
	\R[x] / g(x) \R[x] \to \R \times \R \times \R,
\]
with the isomorphism $[f ]_g \mapsto (f(0), f(1), f(2))$. The inverse map $\R \times \R \times \R \to \R[x] / g(x) \R[x]$ can be calculated via the Lagrange interpolation formula. It is  
\[
	(y_0,y_1,y_2) \in \R^3 \mapsto [ y_0 f_0 + y_1 f_1 + y_2 f_2 ]_g
\]
with 
\begin{align*} 
	f_0(x) & = \frac{1}{2} (x-1) (x-2),
\\ f_1(x) & =  - x (x-2) 
\\ f_2(x) & = \frac{1}{2} x (x-1) 
\end{align*} 
The aim of this example is to highlight that the CRT is a far reaching generalization of the interpolation and, in the case of polynomial rings, it incorporates the classical interpolation tasks. This is due to the fact that for a ring $R$ and an ideal $I$, the canonical homomorphism $R \to R / I$ is the evaluation at the ideal $I$. In the setting with $R = \R[x]$ that we had in the example, such evaluations are indeed some concrete evaluations: for $I = g(x) \R[x]$ this is the evaluation at the set $\{0,1,2\}$ of the roots of $g(x)$ and for $I = (x-a)\R[x]$ it is the evaluation at $a \in \R$. 

 It is actually interesting to take a look at the available interpolation formulas (like Newton interpolation, Neville interpolation and others) and to interpret them in terms of residue class rings. 

\section{Units and inversion in rings}

\subsection{Definition of a unit} In a non-zero ring $R$, the set $R^\times$ is the set  elements $a \in R$ that have a multiplicative inverse in $R$. This is a subset of $R \setminus \{0\}$ that contains $1$. Another common notation for $R^\times$ is $R^\ast$. 

\subsection{Group of units} $(R^\times, \cdot )$ is an Abelian group. 

\begin{proof} 
	If $a, b \in \R^\times$, then $(ab) b^{-1} a^{-1} = 1$ and $b^{-1} a^{-1} (ab) = 1$ so that $ab \in R^\times$, which implies that $\cdot$ can be restricted to $R^\times \times R^\times \to R^\times$. Of course, associativity of $\cdot$ under narrowing the scope to $R^\times$ is inherited from the associativity of $\cdot$ in the whole ring $R$. The $1$ is in $R^\times$ and is a neutral element in $R^\times$. Hence, $R^\times$ is a group. 
\end{proof} 

\subsection{Examples} Here are some easy examples of the groups of units.
\begin{enumerate} 
	\item $k[x_1,\ldots,x_n]^\times = k \setminus \{0\}$
	\item $\Z^\times = \{1,-1\}$. 
	\item $k^\times = k \setminus \{0\}$
\end{enumerate} 

\subsection{Irreducible elements in a ring} It might be useful to have a definition of an irreducible element for an arbitrary non-zero ring $R$.  An element $r \in R$ that is not zero and not a unit is called irreducible if for every decomposition $r = ab$ with $a,b \in R$ one necessarily has $a \in R^\times$ or $b \in R^\times$. According to this definition, an integer $z \in \Z$ is irreducible if and only if $|z|$ is a prime number. The definition also aligns with the notion of irreducibility of polynomials that we introduced for the ring $k[x]$. 

\subsection{Fields} We observe that a non-zero ring $R$ is a field if and only if $R^\times = R \setminus \{0\}$. 

\subsection{An easy lemma for finite Abelian groups} Let $(G, \cdot)$ be a finite Abelian group. Then $a^{|G|} =1$ for every $a \in G$. 

\begin{proof} 
	The map $f : G \to G$, $f(x) = a x$ is bijective as it is invertible with the inverse $f^{-1}(x) = a^{-1} x$. Hence, $\prod_{x \in G} x = \prod_{x \in G} f(x)$. This can be spelled out as 
	\[
		\prod_{x \in G} x = \prod_{x \in G} a x
	\]
	and then rewritten as 
	\[
		\prod_{x \in G} x = \underbrace{ \left( \prod_{x \in G} a \right)}_{a^{|G|}} \left( \prod_{x \in G}  x \right). 
	\]
	So, 
	\[
		p = a^{|G|} p
	\]
	with $p = \prod_{x \in G} x$. Multiplying with $p^{-1}$, we obtain 
	\[
		1 = a^{|G|}. 
	\]
\end{proof} 

\subsection{Reference to Lagrange theorem} The above statement is also true for general finite groups, which are not necessarily Abelian. Still, the proof that we give for the Abelian case is particularly easy and the statement for the Abelian groups will be enough for our purposes. The assertion $a^{|G|} = 1$ in the case of general finite groups $G$ can be derived from the Lagrange theorem, a basic well-known result in group theory. 

\subsection{Corollary for rings} If $R$ is a finite non-zero ring, then for every $a \in R^\times$, one has $a^{|R^\times|} =1$. This can also be formulated as: 
\begin{enumerate} 
\item $a^{|R^\times| + 1} = a$
\item $a^{|R^\times|-1}$ is the inverse element of $a$. 
\end{enumerate} 

\begin{proof} 
	This is true due to the fact that $R^\times$ is an Abelian group. 
\end{proof} 

\subsection{Corollary for fields} If $k$ is a finite field, then $a^{|k|} = a$ for every $a \in k$. In particular, if $a \ne 0$, then  $a^{|k|-1} = 1$ and $a^{|k|-2}$ is the inverse of $a$. 

\begin{proof} 
	For a field $k$, the group of units is $(k \setminus \{0\}, 0)$. So, we can apply the above corollary. 
\end{proof} 

\subsection{Euler's totient function} Totient function, also known as Euler's Totient or Euler phi-function, is a special function $\phi : \N \to \N$. There are several equivalent ways to define it. There are explicit definition on $\phi(n)$ based on the $\gcd$ and the prime factorization of $n$. Such definitions allow one to calculate the totient function, but they have the drawback that they don't clarify the role of this function. 

Following Lang's Algebra (p. 94), for  $n \in \Z_{\ge 2}$ we definite, the $\phi(n) := | (\Z / n \Z)^\times |$. We also define $\phi(1)=1$, which makes sense because in the zero ring $ \Z / \Z$ the unique element (which is both zero and one) is invertible with respect to multiplication. So, conceptually, $\phi(n)$ counts the number of elements in the ring $\Z / n \Z$ that are not invertible. This is what $\phi(n)$ is introduced for.  

Frequently, in the development of a mathematical theory, the conceptual starting point and the later consequences get reversed. A good example is Euler’s totient function. Its original motivation is extremely simple: it just counts the invertible elements of the ring 
$\Z/ n \Z$
 This definition expresses the intention behind the function. Only afterwards do we derive explicit technical formulas. These formulas are not the origin of the concept. They are consequences of studying it.

However, in many textbooks and expositions the order is inverted. The technical formulas are presented first, and the original conceptual motivation is treated as secondary. This reversal can easily create confusion for learners.

\subsection{Invertibility criterion for quotient rings} Let $I$ be an ideal in a ring $R$, then $(R / I )^\times$ is the set of all $[a]_I$, with $a \in R$, where the ideals $a R$ and $I$ are coprime, that means, $a R + I = R$. 
\begin{proof} 
	If $a R + I = R$, then $a u + v=1$ holds for some $u \in R$ and some $v \in I$. But then $[a]_I \cdot [u]_I = 1$, implying that $[a]_I \in (R / I)^\times$. Conversely, if $[a]_I \cdot [u]_I =1$, then $a u + v =1$ for some $v \in I$. But then $a R + I = R$, because the ideal $a R + I$ contains $1$. 
\end{proof} 

\subsection{Maximal ideals} An ideal $I$ is said to be maximal if $I$ is a proper subset of $I$ and there are no ideals lying strictly between $I$ and $R$. 

\subsection{Characterization of maximal ideals in $\Z$ and $k[x]$} 

With the above description, it can be immediately verified that, for $n \in \N$, an ideal $n \Z$ is maximal if and only if $n$ is a prime number. Completely analogously, for $g(x) \in k[x]$, an ideal $g(x) k[x]$ is maximal if and only if $g(x)$ is an irreducible polynomial. 

\subsection{A maximal ideal in $k[x,y]$}
 The ideal $I = \langle x, y\rangle = x k[x,y] + y k[x,y]$ is maximal because it consists of all polynomials that are zero at $(0,0)$. In fact a polynomial being zero at $(0,0)$ has no constant term so that each term of it is divisible by $x$ or $y$. 
 
 If we enlarge this ideal it to  $I + g(x,y) k[x,y]$, where $g(0,0) \ne 0$, then we can write every polynomial $f(x,y)$ as $(f- \frac{f(0,0)}{g(0,0)} g) + \frac{f(0,0)}{g(0,0)} g$. So, $I + g(x,y) k[x,y] = k[x,y]$. Consequently, $I$ is a maximal idea. 
 
\subsection{Maximal ideals and fields} $R/ I$ is a field if and only if the ideal $I$ is maximal. 

\begin{proof}
	Assume that $I$ is maximal. Let us show that every $[f]_I$ with $f \not\in I$ is invertible. We consider the ideal $f R$ and  add it to $I$, we get the ideal $I + f R$ larger than $I$ as it contains $I$ and $f \not\in I$.  By maximality, $I + f R = R$, that is $I$ and $f R$ are coprime. But then, by the above discussion $[f]_I \in R^\times$. So, $R / I$ is indeed a field. 
	
	Conversely, if $I$ is not maximal, then there exists an ideal $J$ that is strictly larger than $I$ but still a strict subset of $R$. Take an element in $J$ but not in $I$. Then $I + f R$ is a subset of $J$ and thus a strict subset of $R$. Hence, the ideals $I$ and $f R$ are not coprime. But this means that $[f]_I \ne 0$ and $[f]_I \not\in R^\times$. It follows that $R / I$ is not a field. 
\end{proof} 

 \subsection{Maximal ideals in $\R[x]$} Interestingly, maximal ideals in $\R[x]$ are in one-to-one correspondence with complexity numbers.  More precisely, every ideal $I\subseteq \R[x]$ can be uniquely written as $I = \{ f \in \R[x] \colon f(c) = 0\}$, where $c \in \C$. Thus, maximal ideals can be viewed as points. Generally, we every ring $R$ one can associate the set $\mathcal{M}$ of all its maximal ideals and consider elements of $R$ as functions over $\mathcal{M}$, for which the evaluation at $I \in \mathcal{M}$ is the operation $f \mapsto [f]_I$. The set of the maximal ideals is called the maximum spectrum of the ring. 


\subsection{Application to the ring $\Z / n \Z$} Remember that for the $\Z$ the condition $m \Z + n \Z = \Z$ is equivalent to $\gcd(m,n)=1$, when $m,n \in \Z$. 

Application of the above to the ring $\Z / n \Z$ gives that 
\[
	(\Z / n \Z)^\times = \{ [z]_n : z \in \{0,\ldots,n-1\}, \gcd(z,n)=1\}.
\]
 In particular, the totient function can be expressed as 
 \[
 	\phi(n) = \Bigl| \{ z \in \{0,\ldots,n-1\} \colon \gcd(z,n)=1\} \Bigr|.
 \]
 We have already discussed when $\Z / n \Z$ is a field. But this discussion, too, 
implies that $ \Z / n \Z$ is a field if and only if $n$ is prime number. 


\subsection{Another property of the Euler totient function} If $n= n_1 \cdots n_s$, where $n_1,\ldots,n_s$ are pairwise coprime, then $\phi(n) = \phi(n_1) \cdots \phi(n_s)$. This follows by application of the CRT and using the fact $(R_1 \times \cdots \times R_s)^\times = R_1^\times \times \cdots \times R_s^\times$, since by the Chinese Remainder theorem, 
\[
	\Z / n_1 \cdots n_s \Z \cong \Z / n_1 \Z \times \cdots \times \Z / n_s \Z. 
\]

\subsection{Calculating the Euler totient function} By the above result, if we have a prime factorization $n = p_1^{m_1} \cdots p_s^{m_s}$ of $n \in \Z_{\ge 2}$ with multiplicities $m_1,\ldots, m_s \in \N$ and pairwise distinct primes $p_1,\ldots,p_s$, we obtian that 
\[
	\phi(n) = \phi(p_1^{m_1}) \cdots \phi(p_s^{m_s}).
\]
So for being able to calculate the totient function $\phi(n)$ out of the prime factorization of $n$, it remains to calculate it on prime powers $p^m$, where $p$ is a prime and $m \in \N$. We need to count the number of invertible elements in $\Z / p^m \Z$. Every number in the range $\{1,\ldots, p^m\}$ that is not divisible by $p$ gives a coset of an invertible element. There are exactly $p^{m-1}$ numbers in the range that are divisible by $p$. Hence, $\phi(p^m) = p^m - p^{m-1} = p^{m-1} (p-1)$. This gives the explicit formula 
\[
	\phi( n) = p_1^{m_1-1} (p_1-1) \cdots p_s^{m_s-1} (p_s-1). 
\]
This formula, known as \emph{Euler's product formula},  is a consequence of the ring theory for the quotient rings of the ring $\Z$ and its proof involves the CRT. 

\section{The RSA cryptosystem from the abstract perspective} 

Those of you who are in the cybersecurity program and will take cryptography later will see RSA again. Here, I just want to point out that the key ideas behind RSA can already be understood from a basic ring-theoretic perspective, without doing the detailed calculations. Of course, the calculations matter in practice, but they are not the reason why the encryption and decryption are inverse functions and why decryption really returns the plain text. That structural reason is what students often miss. Not every calculation explains why something works. Understanding the structure is crucial. 

To make the  structure more visible, we refrain from using the concrete ring $R =\Z$ on which the RSA is based. The RSA computations in $\Z$ are done using the EEA, but when we move over to a general ring $R$, we deliberately switch that part off which forces us to focus on the structural understanding. 

\subsection{The ring $R / ( I \cap J)$ for distinct maximal ideals $I$ and $J$.} If $I$ and $J$ are maximal ideals with $I \ne J$ then $I +J = R$.
Because $I+J$ is larger than both $I$ and $J$ maximal, $I+J$ is necessarily $R$. 

Assume now that$ R / I$ and $R/ J$ are finite of size $p$ and $q$, respectively. Then, by the CRT $a \in R / (I \cap J)$ is finite of size $pq$, because $R / (I \cap J) \cong R / I \times R / J$. 

As $I$ and $J$ are maximal, both $R / I$ and $R/ J$ are fields. So, $R / (I \cap J)$ is isomorphic to the product of two fields. In the field $R / I$ we have the identity $x^{p-1} =1$ for every $x \in \R / I$ with $x \ne 0$. In the field $R /J$ we have the identity $y^{q-1} = 1$ for every $y \in R / J$ with $y \ne 0$. We can raise those identities to any power. So, in particular, we have $x^{m (p-1)(q-1)} = 1$ and $y^{m (p-1)(q-1)} =1$ valid for each $m \in \N$. Multiplying by $x^e$ with $e \in \N$  we make them into identities 
\begin{align*} 
	x^{e + m (p-1) (q-1)} & = x^e &  y^{e+ m (p-1) (q-1)} & = y^e
\end{align*} 
valid for all elements $x \in R/I$ and $y \in R/J$. 
But this means that $z^{e+ m (p-1) (q-1)} = z^e$ holds for $z \in R / I \times R / J$, which implies that 
\[
	a^{e + m (p-1)(q-1)} = a^e
\] holds for the $a \in R / (I \cap J)$. 

\subsection{Power functions over $R /( I \cap J)$} In the above setting, let us consider the power functions $\pow_e : R / (I \cap J) \to R  (I \cap J)$, $\pow_e(x):= x^e$. The above observations show that letting $e$ go through the range $e \in \{1,\ldots,(p-1)(q-1) \}$ we obtain all power functions, since the power functions are repeated periodically as the power functions differing by a multiple of $(p-1)(q-1)$ coincide. 
In particular, that means that if two power functions $e,d \in \N$ are such that $e d = m (p-1) (q-1)$ then they are inverse to each other in the sense of the equality $\pow_e(\pow_d(x))=\pow_d(\pow_e(x))=x$. 


\subsection{Rivest-Shamir-Adleman cryptosystem} This cryptosystem is abbreviated as RSA. The RSA cryptosystem uses $ R = \Z$, $I = p \Z$ and $J = q \Z$, where $p$ and $q$ are distinct prime numbers, so that $I \cap J = p q \Z$. The public key is the power $(e, p q )$ with $e \in \N$ and $pq \in \N$. The public key determines the power map $\pow_e(x) = x^e$ on the ring $ \Z / p q \Z$. 
So, it  allows the sender to send a plain text $x \in \Z / p q \Z$ as $x^e \in \Z / pq \Z$. The receiver, holding the private key  $d \in \N$, which satisfies $e d \equiv 1 \mod (p-1)(q-1)$, inverts the power function $\pow_e(x)$  by applying the power function $\pow_d$. 

For being able to break the RSA cryptosystem, the interceptor needs to calculate private key $d$.If one gets the factor $p$ and $q$ from $pq$, then one can determine $(p-1)(q-1)$ and calculate $d$ using $e$ and $(p-1)(q-1)$ by means of the EEA. However, no efficient algorithm for factorization of integers into prime factors is available for the standard models of computations. Theoretically, quantum computers can factorize integers in polynomial time, but so far large scale quantum computers are not available. Thus, RSA is relying on hardness of the factorization of integers. 



\end{document} 